\documentclass[11pt]{article}

% =====================================================================
% PEARL — Empirical Findings on Internal vs External Generalization
% =====================================================================
% This document is a self-contained analysis section. Drop it into
% Overleaf, compile, and the figures load from ../analyhsis/figures/.
%
% All metrics quoted below are reproducible from:
%   results/ablation/sweep_matrix.csv        (internal)
%   results/external_validation/summary.csv  (external with preproc)
%   results/external_validation/comparison_with_vs_noproc.csv
%   results/external_validation/comparison_with_vs_noproc.csv (noproc)
% =====================================================================

\usepackage[margin=1in]{geometry}
\usepackage{graphicx}
\usepackage{booktabs}
\usepackage{longtable}
\usepackage{amsmath}
\usepackage{hyperref}
\usepackage{xcolor}
\usepackage{caption}
\usepackage{subcaption}
\usepackage{float}

\hypersetup{colorlinks=true, linkcolor=blue!50!black, urlcolor=blue!50!black}
\graphicspath{{../analyhsis/figures/}}

\title{PEARL — Cross-dataset Generalization Analysis\\[0.2em]
\large Internal saturation vs external collapse on Figshare~$\rightarrow$~PCOSgen}
\author{Empirical findings from the 18-run ablation}
\date{\today}

\begin{document}
\maketitle

% =====================================================================
\section{Summary}
% =====================================================================

We trained 9 architectures (ResNet-50/101, DenseNet-121/169, EfficientNet-B0,
ConvNeXt-Tiny, MobileNetV3-Large, ViT-B/16, Swin-Tiny) under 2 denoising
preprocessing pipelines (SRAD, Gaussian) on the Figshare PCOS dataset, totalling
18 transfer-learning runs. We then evaluated each run on the held-out PCOSgen
cohort (Sundari et al.\ 2025) to measure cross-dataset generalization.

Three findings dominate:

\begin{enumerate}
  \item All 18 configurations saturate on the internal Figshare test set
        (AUC $\geq 0.999$), making the internal ablation uninformative for
        architecture selection.
  \item All 18 configurations collapse to near-random performance on the
        external PCOSgen cohort (AUC 0.45--0.59, balanced accuracy $\sim 0.50$).
        The drop is structural and uniformly distributed across architectures.
  \item The collapse is \emph{not} caused by the SRAD/Gaussian preprocessing:
        disabling preprocessing yields essentially identical external AUC
        (mean $\Delta = +0.022$ AUC in preprocessing's favor). The
        discriminative features simply do not transfer between the two datasets.
\end{enumerate}

The remainder of this document presents the evidence for each finding in
detail, with all 18 (prep $\times$ arch) configurations traced end-to-end.

% =====================================================================
\section{Data context}
% =====================================================================

\subsection{Internal: Figshare PCOS}

\begin{itemize}
  \item Source: \texttt{data/\{infected,noninfected\}} (3,996 unique images
        after md5 dedup; 3,184 PCOS / 812 non-PCOS).
  \item Stratified 80/10/10 split, seed 42, group-aware on filename stem
        (\texttt{src/data/splitter.py}).
  \item Class imbalance: 3.9:1 PCOS:non-PCOS.
  \item Preprocessed with SRAD (\emph{Speckle-Reducing Anisotropic Diffusion},
        ultrasound-tuned) and Gaussian denoising pipelines.
\end{itemize}

\begin{figure}[H]
  \centering
  \includegraphics[width=0.85\linewidth]{fig11_internal_split_composition.png}
  \caption{Internal Figshare split composition. The test set has 81 negatives,
           so a single correct/incorrect prediction swings AUC by $\sim0.001$.
           This makes the internal test set underpowered for architecture
           discrimination.}
  \label{fig:internal-split}
\end{figure}

\subsection{External: PCOSgen}

\begin{itemize}
  \item Source: \texttt{data\_external/pcosgen/} (Sundari et al.\ 2025,
        Kaggle upload).
  \item Cohort: 4,668 images (3,627 PCOS / 1,041 healthy). 3.5:1 imbalance.
  \item Ground-truth labels read from the dataset's own
        \texttt{class\_label.csv} (\texttt{Healthy} column, 1=healthy/0=PCOS),
        cross-referenced against folder names (\texttt{infected/, healthy/}).
        The two encodings \textbf{agree} (see~\ref{fig:csv-confusion}).
\end{itemize}

\begin{figure}[H]
  \centering
  \includegraphics[width=0.65\linewidth]{fig12_external_split_composition.png}
  \caption{PCOSgen external cohort. 4,668 images, 71.8\% prior on PCOS,
           raw images 300$\times$300 vs Figshare's 850$\times$578.}
  \label{fig:external-split}
\end{figure}

\subsection{Image acquisition differs}

\begin{figure}[H]
  \centering
  \includegraphics[width=0.7\linewidth]{fig13_image_resolution.png}
  \caption{Raw image resolution. Figshare median $\sim$850$\times$578, PCOSgen
           median 300$\times$300. The size difference is a proxy for
           different scanner / acquisition settings.}
  \label{fig:resolution}
\end{figure}

% =====================================================================
\section{Finding 1: Internal test set is saturated}
% =====================================================================

The internal figshare-test AUC is $\geq 0.999$ for every (prep $\times$ arch)
combination. The spread across the 18 runs is only 0.0008 AUC.

\begin{figure}[H]
  \centering
  \includegraphics[width=\linewidth]{fig1_internal_auc.png}
  \caption{Internal test AUC. Every run is at or above 0.999. The ablation
           cannot rank architectures on this test set.}
  \label{fig:internal-auc}
\end{figure}

The full per-run breakdown is in \texttt{results/ablation/sweep\_matrix.csv}.
The single best and worst configurations differ by 0.0005 AUC and 0.022 MCC,
neither of which is statistically meaningful at $n=406$ with 81 negatives.

% =====================================================================
\section{Finding 2: External validation collapses to near-random}
% =====================================================================

\subsection{The AUC drop is uniform across all 18 configurations}

\begin{figure}[H]
  \centering
  \includegraphics[width=\linewidth]{fig5_internal_vs_external_paired.png}
  \caption{Internal vs external AUC. Every (prep $\times$ arch) drops from
           $\geq 0.999$ to 0.45--0.59. The drop is uniformly distributed across
           architectures and preprocessing pipelines.}
  \label{fig:internal-vs-external}
\end{figure}

\begin{figure}[H]
  \centering
  \includegraphics[width=0.85\linewidth]{fig6_drop_distribution.png}
  \caption{Distribution of internal$-$external AUC drops across all 18 runs.
           Mean $\Delta = 0.464$, median $\Delta = 0.479$, minimum $\Delta = 0.412$.
           \textbf{No configuration escapes the collapse.}}
  \label{fig:drop-dist}
\end{figure}

\subsection{External AUC rank}

\begin{figure}[H]
  \centering
  \includegraphics[width=0.95\linewidth]{fig14_external_auc_rank.png}
  \caption{External PCOSgen AUC, sorted best-to-worst. Best: srad/ResNet-50
           at 0.587. Worst: srad/DenseNet-169 at 0.450. Spread 0.137 AUC,
           which is within the dataset's $\sim 0.04$ CI width at N=4668
           with 71.8\% prior.}
  \label{fig:ext-rank}
\end{figure}

\subsection{The collapse is a "predict everything as PCOS" failure}

\begin{figure}[H]
  \centering
  \includegraphics[width=\linewidth]{fig3_external_specificity.png}
  \caption{External specificity. Every model correctly identifies fewer than
           7.3\% of true healthy PCOSgen cases. The model has learned a
           class prior, not a discriminative boundary.}
  \label{fig:spec}
\end{figure}

\begin{figure}[H]
  \centering
  \includegraphics[width=0.85\linewidth]{fig4_external_mcc.png}
  \caption{External MCC. The Matthews Correlation Coefficient, which is
           insensitive to class imbalance, is $\leq 0.04$ for every
           configuration. This is the single number that best summarises
           the failure.}
  \label{fig:mcc}
\end{figure}

\begin{figure}[H]
  \centering
  \includegraphics[width=0.7\linewidth]{fig9_sens_vs_fpr.png}
  \caption{Operating point on the ROC plane (False-Positive Rate vs
           True-Positive Rate). Every model is clustered in the top-right
           corner: high sensitivity, very high false-positive rate. The
           classifiers are not finding a decision boundary that separates
           the two classes on PCOSgen.}
  \label{fig:roc-op}
\end{figure}

\subsection{All-metric heatmap}

\begin{figure}[H]
  \centering
  \includegraphics[width=\linewidth]{fig18_metric_heatmap.png}
  \caption{All six headline metrics on PCOSgen. AUC is in the yellow band
           (0.45--0.59, no signal), F1/Accuracy are green only because the
           dataset is 71.8\% PCOS (majority-class baseline), and MCC is
           uniformly red ($\leq 0.04$). Only the right column honestly
           summarises the failure.}
  \label{fig:heatmap}
\end{figure}

\subsection{Class-prior operating point dominates}

The dataset's PCOS prior is 71.8\%. The "always predict PCOS" baseline
achieves accuracy 0.718, sensitivity 1.0, specificity 0.0, F1 $\approx 0.84$,
MCC 0.0, AUC 0.5. Most of our models are within a few points of this
degenerate baseline. The single metric that breaks the degeneracy is
specificity (and the derived MCC), which is why those are the headline
external numbers in this document.

% =====================================================================
\section{Finding 3: Label polarity is correct (not flipped)}
% =====================================================================

A legitimate concern is that PCOSgen's CSV uses different label conventions
than the internal dataset. We cross-referenced 3,200 per-image predictions
against the dataset's own \texttt{class\_label.csv}:

\begin{figure}[H]
  \centering
  \includegraphics[width=0.7\linewidth]{fig17_confusion_matrix_csv_check.png}
  \caption{Confusion matrix computed using PCOSgen's CSV as ground truth
           (srad/ResNet-50 evaluated on the 3,200 images that match between
           folder names and CSV rows). TP=2228, FP=862, FN=69, TN=41.
           The model genuinely over-predicts PCOS (sensitivity 0.97,
           specificity 0.045). The labels are \textbf{not} flipped.}
  \label{fig:csv-confusion}
\end{figure}

\textbf{Conclusion:} the external collapse is not a labelling bug. The
folder convention \texttt{infected}$\leftrightarrow$\texttt{Healthy=0}
correctly matches the CSV's encoding ($1$=healthy, $0$=PCOS per
\texttt{info.txt}).

% =====================================================================
\section{Finding 4: Preprocessing is not destroying the images}
% =====================================================================

A second concern is that the SRAD / Gaussian pipelines tuned on Figshare
might be inappropriate for PCOSgen's acquisition characteristics. We
visually inspected the preprocessed outputs on both datasets:

\begin{figure}[H]
  \centering
  \includegraphics[width=\linewidth]{fig15_visual_sanity_check.png}
  \caption{Visual sanity check. Three rows (Figshare/infected, PCOSgen/healthy,
           PCOSgen/infected), four columns (raw, SRAD-preprocessed,
           Gaussian-preprocessed, no-preprocessing). All four modes preserve
           the ultrasound fan-shaped anatomy on both datasets. The
           preprocessor is functionally correct for PCOSgen images.}
  \label{fig:visual}
\end{figure}

The preprocessed outputs were originally rendered as black (float32
normalized arrays clipped to uint8 without stretching). After proper
1st--99th-percentile stretching, the anatomical structure is intact on
both datasets. The preprocessor is not the culprit.

% =====================================================================
\section{Finding 5: With vs without preprocessing — same collapse}
% =====================================================================

To test whether preprocessing is what couples the model to the wrong
intensity distribution, we re-ran all 18 external evaluations with
preprocessing \textbf{disabled} (resize + ImageNet normalize only). The
outputs land in \texttt{<run\_dir>/external\_validation\_noproc/} and
do not overwrite the existing with-preprocessing results.

\begin{figure}[H]
  \centering
  \includegraphics[width=0.85\linewidth]{fig16_external_noproc_auc.png}
  \caption{External AUC with preprocessing disabled. The same collapse
           pattern: AUC 0.36--0.58, no architecture exceeds 0.59.}
  \label{fig:noproc-auc}
\end{figure}

\begin{figure}[H]
  \centering
  \includegraphics[width=0.85\linewidth]{fig7_with_vs_noproc_scatter.png}
  \caption{Scatter of with-preproc vs no-preproc external AUC. Most points
           are clustered between 0.40 and 0.58 on both axes. The diagonal
           trend shows that preprocessing is a marginal effect, not a
           categorical lever.}
  \label{fig:scatter}
\end{figure}

\begin{figure}[H]
  \centering
  \includegraphics[width=\linewidth]{fig8_delta_preprocessing.png}
  \caption{Per-architecture $\Delta$ AUC (with-preproc $-$ no-preproc).
           Mean $\Delta = +0.022$ AUC in preprocessing's favor. The spread
           is large ($-0.06$ to $+0.15$), so preprocessing neither helps
           nor hurts consistently.}
  \label{fig:delta}
\end{figure}

\subsection{What turning off preprocessing changes}

Three runs gain $\Delta \geq +0.06$ AUC from preprocessing (the resnet101,
resnet50, densenet121 cases). Five runs lose $\leq 0.06$ AUC from
preprocessing (convnext\_tiny, vit\_base, mobilenetv3\_large,
densenet169). The cross-architecture spread is larger than the
preprocessing effect, meaning \textbf{architecture choice and preprocessing
choice are both noise at the external evaluation scale}.

% =====================================================================
\section{Finding 6: Balanced accuracy is the cleanest summary}
% =====================================================================

The balanced accuracy (mean of sensitivity and specificity) is the
single number that best captures the model's true discriminative ability
on a class-imbalanced external set.

\begin{figure}[H]
  \centering
  \includegraphics[width=\linewidth]{fig10_balanced_accuracy.png}
  \caption{Balanced accuracy: internal $\approx 0.99$ vs external $\approx 0.50$.
           Every classification threshold on PCOSgen produces a balanced
           accuracy close to random (0.5).}
  \label{fig:bal-acc}
\end{figure}

% =====================================================================
\section{Reproducibility}
% =====================================================================

All artifacts are committed to the repository:

\begin{itemize}
  \item \texttt{results/ablation/sweep\_matrix.csv} --- all 18 runs' internal
        figshare-test metrics.
  \item \texttt{results/external\_validation/summary.csv} --- aggregated
        external (PCOSgen) metrics with original SRAD/Gauss preprocessing.
  \item \texttt{results/external\_validation/comparison\_with\_vs\_noproc.csv}
        --- side-by-side comparison with preprocessing disabled.
  \item \texttt{results/ablation/checkpoints/<prep>/<arch>/external\_validation/pcosgen.json}
        --- per-arch per-prep JSON metrics with preproc.
  \item \texttt{results/ablation/checkpoints/<prep>/<arch>/external\_validation\_noproc/pcosgen\_noproc.json}
        --- per-arch per-prep JSON metrics without preproc.
  \item \texttt{results/probe\_preproc/} --- visual sanity check artefacts
        (raw side-by-side images, preprocessed arrays, percentile-stretched
        renderings).
  \item \texttt{results/external\_validation/comparison\_with\_vs\_noproc.csv}
        --- tidy CSV with all 18 rows, both preproc modes.
\end{itemize}

Scripts (in \texttt{scripts/}) that produce the above:

\begin{itemize}
  \item \texttt{eval\_external\_all.py} --- fanout driver for with-preproc eval.
  \item \texttt{eval\_external.py} --- single-run with-preproc evaluator.
  \item \texttt{eval\_external\_noproc.py} --- single-run no-preproc evaluator.
  \item \texttt{eval\_external\_noproc\_all.py} --- fanout driver for no-preproc.
  \item \texttt{generate\_analysis\_figures.py} --- regenerates every figure
        in this document.
\end{itemize}

Total external inference cost: 36 (18 with preproc + 18 without) $\times$
$\sim$70s/run on a single RTX 4070.

\end{document}
